% !PRINTED 254
\weberpage{254}
But by (12) $\beta$ is divisible both by $e$ and by $e_2$, and
$\beta e e_2$ is a square; hence
\[
\left(\frac{\alpha}{\beta}\right) = \left(\frac{\alpha}{ee_2}\right) = \left(\frac{\alpha}{e}\right)\left(\frac{\alpha}{e_2}\right) = \left(\frac{\delta}{e}\right)\left(\frac{\alpha}{e_2}\right) \ (\text{since}\ \alpha\delta \equiv 1\ \mathrm{mod}\ e);
\]
furthermore, by (12),
\[
\left(\frac{\delta}{e}\right) = \left(\frac{c}{e}\right)\left(\frac{\partial_2}{e}\right), \quad \left(\frac{\alpha}{e_2}\right) = \left(\frac{c_2}{e_2}\right)\left(\frac{-\partial}{e_2}\right),
\]
and
\[
\left(\frac{\partial_2}{e}\right) = \left(\frac{e_2}{e}\right), \quad \left(\frac{-\partial}{e_2}\right) = \left(\frac{-e}{e_2}\right), \quad \left(\frac{e_2}{e}\right)\left(\frac{-e}{e_2}\right) = (-1)^{\frac{(e+1)(e_2-1)}{4}},
\]
hence
\[
\left(\frac{\alpha}{\beta}\right) = (-1)^{\frac{(e+1)(e_2-1)}{4}}\left(\frac{c}{e}\right)\left(\frac{c_2}{e_2}\right)
\]
\[
i^{\frac{1-\beta}{2}} = i^{\frac{1-ee_2}{2}}; \quad i^{\frac{a-a_2}{2}} = i^{\frac{e-e_2}{2}}.
\]

Thus, under the transformations (2),
\weq{14}{P_{c,\partial,a}\ \text{is carried to}\ e^{\frac{2\pi i}{3}\frac{\lambda-1}{8}}P_{c_1,\partial,a}\ \text{and to}\ e^{\frac{2\pi i}{3}\frac{\beta(\alpha+\delta)-(\beta^2-1)\alpha\gamma}{8}}P_{c_2,\partial_2,a_2}.}

If $n$ is moreover not divisible by 3, then let $c$ be assumed
divisible by 3, whereby
\[
\lambda \equiv 1, \quad \alpha \equiv \delta \equiv 0\ (\mathrm{mod}\ 3)
\]
result, and the third roots of unity occurring in (14) receive the
value 1.

From this we obtain the following theorems:

5. \gesperrt{If $n$ is an odd square number, $c$ divisible by 8, then
the quantities}
\[
P^3_{c,\partial,a},
\]
\gesperrt{and if $n$ is an odd square number not divisible by 3, $c$
divisible by 24, then the quantities}
\[
P_{c,\partial,a}
\]
\gesperrt{are roots of multiplier equations of the first level.}

The first examples are $n=9$, $n=25$.

These equations can be formed in various ways. Here, following
\gesperrt{Kiepert}$^1$), we take the following course: by § 69, we express
the roots of a transformation equation belonging to degree $p$
%CONT>>

\begingroup
\renewcommand{\thefootnote}{1)}
\footnotetext{Zur Transformationstheorie der elliptischen
Funktionen. Crelles Journal, Bd. 87, 88, 95.}
\endgroup
